\begin{abstract}
Vision Transformers scale quadratically with input resolution, making high-resolution training prohibitively expensive, yet detection and segmentation demand it. A natural alternative is resolution extrapolation: train at low resolution and deploy at higher resolution without fine-tuning. Whether this works depends entirely on the positional encoding. Modern encodings such as RoPE, ALiBi, YaRN, and FIRE enable striking length extrapolation in language models, but their behavior on 2D image grids is largely unknown. We present a systematic study, training a single ViT-Tiny on ImageNet-100 at $224\!\times\!224$ and evaluating zero-shot up to $1024\!\times\!1024$ ($4.57\times$) on classification, COCO detection, and ADE20K segmentation. Across all three tasks, additive attention-bias methods consistently outperform rotation-based methods at large scales: at $4.57\times$, FIRE retains $63.1\%$ Top-1 accuracy while RoPE collapses to $18.9\%$. An attention-entropy analysis shows that bias-based encodings preserve focused, semantically coherent attention, whereas rotation-based phases drift out of distribution and induce attention collapse. These results recast extrapolation robustness as the primary axis for choosing positional encodings, and yield a practical recipe for resolution-flexible Vision Transformers.
%\keywords{Vision Transformers \and Positional Embeddings \and Resolution Extrapolation \and RoPE \and ALiBi \and FIRE}
\end{abstract}
